\chapter{Introduction}

The interaction between superconductivity and ferromagnetism has been a topic of interest in condensed matter physics due to the competing nature of these two phenomena. Superconductors exhibit zero electrical resistance and expel magnetic fields (Meissner effect), while ferromagnets have a non-zero magnetization that generates internal magnetic fields. Jeon et al. suggested that in superconducting-ferromagnetic-superconducting (SC-FM-SC) heterostructures, the presence of a ferromagnetic layer can induce spin-triplet Cooper pairs \cite{Jeon2018}, having observed an increase in the Gilbert damping below $T_\textrm{C}$. However, the study conducted ferromagnetic resonance (FMR) measurements only at one frequency ($f = \SI{20}{GHz}$) and did not investigate the frequency dependence of the Gilbert damping parameter $\alpha$. Müller et al. \cite{muller2021temperature} and Mangold (unpublished) conducted broadband FMR measurements on SC-FM-SC heterostructures, but did not find evidence for spin-triplet supercurrents. Mangold observed unusual spikes in the Gilbert damping and other parameters near the critical temperature $T_\textrm{C}$ of the superconductor, indicating a potential interplay between superconductivity and magnetization dynamics. This spiking behavior has not been reported in literature yet, and motivated the investigation in this thesis to understand the underlying physics.

The goal of this thesis is twofold: First, to investigate the Gilbert damping constant $\alpha$, the inhomogeneous broadening $\Delta H_0$, the effective magnetization $M_\textrm{eff}$, and the resonance field shift $H_\textrm{shift}$ as functions of temperature in SC-FM-SC heterostructures near the critical temperature $T_\textrm{C}$ using broadband FMR for Nb thicknesses of $d_\textrm{Nb} = \SI{45}{nm}$ and $d_\textrm{Nb} = \SI{30}{nm}$, and second, to fabricate a device capable of FMR measurements under a DC-bias current. The investigation will focus on understanding the influence of superconductivity on the magnetization dynamics of the ferromagnetic layer, particularly near $T_\textrm{C}$. The fabrication of a device capable of FMR measurements under a DC-bias current will enable further studies on the interplay between superconductivity and magnetization dynamics in SC-FM-SC heterostructures.
